\section{Элементарные приложения}

\subsection{Неравенство обработки данных и неравенство Фано}

В данном подразделе мы рассматриваем исключительно дискретные случайные величины. Основание логарифма подразумевается равным 2. 

\begin{definition}[Марковская цепь]
	Случайные величины $\xi, \eta, \zeta$ со значениями в $\mathcal{X}, \mathcal{Y},\mathcal{Z}$ соответственно образуют марковскую цепь (обозначение: $\xi \to \eta \to \zeta$), если для всех $x \in \mathcal{X}, y \in \mathcal{Y}, z \in \mathcal{Z}$, таких что $\P(\eta = y, \zeta = z) > 0$,
	\[
		\P(\zeta = z | \eta = y, \xi = x) = \P(\zeta = z | \eta = y).
	\]
\end{definition}
\begin{remark}
	В терминах функций вероятности, есть несколько эквивалентных формулировок
	$p_{\xi\eta\zeta}(x,y,z) = p_\xi(x) p_{\eta|\xi}(y|x) p_{\zeta|\eta}(z|y)$,
	или, что то же самое, $p_{\zeta|\xi\eta}(z|x,y) = p_{\zeta|\eta}(z|y)$.
\end{remark}
\begin{remark}
	если при условии $\eta$ величины $\xi$ и $\zeta$ независимы. 
\end{remark}
\begin{lemma}[Условная независимость в марковской цепи]
	\label{lem:cond-indep-markov}
	Если $\xi \to \eta \to \zeta$, то условная взаимная информация $\I(\xi; \zeta | \eta) = 0$.
\end{lemma}
\begin{proof}
	По определению условной взаимной информации:
	\[
		\I(\xi; \zeta | \eta) = \sum_{x,y,z} p_{\xi\eta\zeta}(x,y,z) \log_m \frac{p_{\xi\zeta|\eta}(x,z|y)}{p_{\xi|\eta}(x|y) p_{\zeta|\eta}(z|y)}.
	\]
	Из марковского свойства $p_{\xi\zeta|\eta}(x,z|y) = p_{\xi|\eta}(x|y) p_{\zeta|\eta}(z|y)$, поэтому логарифм равен нулю для всех слагаемых, и вся сумма обращается в ноль.
\end{proof}

\begin{theorem}[Неравенство обработки данных (Data Processing Inequality)]
	\label{thm:data-processing}
	Пусть $\xi \to \eta \to \zeta$ образуют марковскую цепь. Тогда:
	\[
		\I(\xi; \eta) \geq \I(\xi; \zeta).
	\]
	Иными словами, никакая обработка данных $\eta$ не может увеличить количество информации о $\xi$.
\end{theorem}
\begin{proof}
	Рассмотрим величину $\I(\xi; \eta; \zeta)$ — взаимную информацию трёх величин. С одной стороны:
	\[
		\I(\xi; \eta; \zeta) = \I(\xi; \zeta) - \I(\xi; \zeta | \eta).
	\]
	С другой стороны:
	\[
		\I(\xi; \eta; \zeta) = \I(\xi; \eta) - \I(\xi; \eta | \zeta).
	\]
	Из леммы \ref{lem:cond-indep-markov} следует, что $\I(\xi; \zeta | \eta) = 0$. Подставляя это в первое равенство, получаем:
	\[
		\I(\xi; \eta; \zeta) = \I(\xi; \zeta).
	\]
	Из второго равенства имеем $\I(\xi; \eta; \zeta) = \I(\xi; \eta) - \I(\xi; \eta | \zeta) \leq \I(\xi; \eta)$, поскольку условная взаимная информация неотрицательна. Следовательно, $\I(\xi; \zeta) \leq \I(\xi; \eta)$, что и требовалось доказать.
\end{proof}

\begin{theorem}[Неравенство Фано]
	\label{thm:fano}
	Пусть $\eta \to \xi \to \hat{\eta}$ — марковская цепь, причём $\eta$ и $\hat{\eta}$ принимают значения в одном и том же конечном множестве $\mathcal{K}$. Обозначим вероятность ошибки:
	\[
		P_{\mathrm{error}} = \P(\eta \neq \hat{\eta}).
	\]
	Тогда справедлива следующая оценка:
	\[
		\H_m(\eta | \xi) \leq \H_m(P_{\mathrm{error}}) + P_{\mathrm{error}} \cdot \log_m(|\mathcal{K}| - 1),
	\]
	где $\H_m(P_{\mathrm{error}})$ — энтропия бернуллиевской случайной величины с вероятностью успеха $P_{\mathrm{error}}$, то есть:
	\[
		\H_m(P_{\mathrm{error}}) = -P_{\mathrm{error}} \log_m P_{\mathrm{error}} - (1-P_{\mathrm{error}}) \log_m (1-P_{\mathrm{error}}).
	\]
\end{theorem}
\begin{proof}
	Применим неравенство обработки данных (теорема \ref{thm:data-processing}) к цепи $\eta \to \xi \to \hat{\eta}$. Имеем:
	\[
		\I(\eta; \xi) \geq \I(\eta; \hat{\eta}).
	\]
	Заметим что $\I(\eta; \xi) = \H(\eta) - \H(\eta | \xi)$. $\I(\eta; \hat{\eta}) = \H(\eta) - \H(\eta | \hat{\eta})$. Откуда:
	\[
		\H(\eta | \xi) \leq \H(\eta | \hat{\eta})
	\]
	Теперь оценим $\H(\eta | \hat{\eta})$. Рассмотрим случайную величину $E = \ind\set{\{\eta \neq \hat{\eta}\}}$ --- индикатор ошибки. Применим правило цепочки для условной энтропии:
	\begin{align*}
		\H(\eta, E | \hat{\eta}) &= \H(\eta | \hat{\eta}) + \underbrace{\H(E | \eta, \hat{\eta})}_{0} \\
		\H(\eta, E | \hat{\eta}) &= \H(E | \hat{\eta}) + \H(\eta | E, \hat{\eta})
	\end{align*}
	Приравнивая правые части, получаем:
	\[
		\H(\eta | \hat{\eta}) = \H(E | \hat{\eta}) + \H(\eta | E, \hat{\eta}).
	\]
	Оценим каждое слагаемое. Имеем $\H(E | \hat{\eta}) \leq \H(E) = \H(P_{\mathrm{error}})$, так как условная энтропия не превосходит безусловную. Распишем второе слагаемое:
	\begin{multline*}
		\H(\eta| E, \hat{\eta}) = \left[
		\sum_{\hat{y} \in \mathcal{K}}\P(\eta = \hat{\eta}, \hat{\eta} = \hat{y}) \sum_{y \in \mathcal{K}}\overbrace{\P(\eta = y | \eta = \hat{\eta}, \hat{\eta} = \hat{y})\log\P(\eta = y | \eta = \hat{\eta}, \hat{\eta} = \hat{y})}^{0}\right. +\\+
	\left.\sum_{\hat{y} \in \mathcal{K}}\P(\eta \neq \hat{\eta}, \hat{\eta} = \hat{y}) \sum_{y \in \mathcal{K}}\P(\eta = y | \eta \neq \hat{\eta}, \hat{\eta} = \hat{y})\log\P(\eta = y | \eta \neq \hat{\eta}, \hat{\eta} = \hat{y})\right]
	=\\= \sum_{\hat{y} \in \mathcal{K}} \P(\hat{\eta} = y, \eta \neq y) \sum_{\substack{y \in \mathcal{K} \\ y \neq \hat{y}}} \P(\eta = y | \eta \neq \hat{y}, \hat{\eta} = \hat{y}) \log \P(\eta = y | \eta \neq \hat{y}, \hat{\eta} = \hat{y}) \leq\\\leq
		\sum_{\hat{y} \in \mathcal{K}} \P(\hat{\eta} = \hat{y}, \eta \neq \hat{y}) \cdot \log(\abs{\mathcal{K}} - 1) =  P_{\mathrm{error}} \cdot \log(|\mathcal{K}| - 1)
	\end{multline*}
	Итак имеем:
	\[
		\H(\eta | \xi) \leq \H(\eta | \hat{\eta}) = \H(E | \hat{\eta}) + \H(\eta | E, \hat{\eta}) \leq \H(E) +  P_{\mathrm{error}} \cdot \log(|\mathcal{K}| - 1)
	\]
	Что в точности есть заключение теоремы.
\end{proof}

\begin{corollary}[Принцип \enquote{Garbage In --- Garbage Out}]
	Пусть по наблюдаемым признакам $\xi$ строится оценка некоторой ненаблюдаемой величины $\eta$: $\hat{\eta} = f(\xi)$; чем меньше $\I(\eta; \xi)$, тем хуже будет оценка. Иными словами \textbf{мусор на входе --- мусор на выходе}
\end{corollary}
\begin{corollary}
	В задаче бинарной классификации, вероятность ошибки оценивается как
	\[
		\P(\text{\normalfont predicted label } \neq \text{\normalfont  true label}) > \H(\text{\normalfont true label } | \text{\normalfont  observed data}) 
	\]
\end{corollary}
\subsection{Отбор признаков на основе взаимной информации}
Итак, мы понимамем, что чтобы решать задачу предсказания, необходимо максимизировать взаимную информацию между целевой переменной $\eta$ и имеющимися признаками $\xi = (\xi_1, \ldots, \xi_p)$. С точки зрения теории, это означает что нужно иметь как можно больше признаков: по правилу цепочки
\[
	\I(\xi; \eta) = \sum_{i = 1}^{p}\I(\xi_i; \eta | \xi_{i - 1}, \ldots, \xi_1)
\]
взаимная информация будет увелиичиваться (не уменьшаться) с ростом $p$; На практике однако есть две проблемы.
\begin{itemize}[compact]
	\item Производительность: огромное число признаков означает большее время обучения для большинства моделей; к тому же многие классические методы теряют численную устойчивость в больших размерностях. 
	\item Точность: многие модели можно обмануть всевозможными странными зависимостями, присутствующими в данных.
\end{itemize}
Инженер по данным должен сделать сигнал максимально легко идентифицируемым для модели, чтобы она могла его изучить. На практике это означает, что \textbf{отбор признаков} может являться важным этапом предобработки. Отбор признаков помогает сконцентрироваться на релевантных переменных в наборе данных, а также может помочь устранить коллинеарные переменные. Он помогает уменьшить шум в наборе данных и помогает модели улавливать релевантные сигналы.

\subsubsection{Методы фильтрации}
В описанной ситуации мы обычно имеем матрицу данных высокой размерности $\textbf{X} \in \mathbb{R}^{n \times p}$ и целевую переменную $y$ (дискретную или непрерывную). Алгоритм отбора признаков выберет подмножество из $k \ll p$ столбцов, $\textbf{X}_S \in \mathbb{R}^{n \times k}$, которые наиболее релевантны целевой переменной $y$.

В общем случае алгоритмы отбора признаков можно разделить на три класса:
\begin{enumerate}[compact]
	\item \textbf{Методы-обёртки} используют обучающиеся алгоритмы на исходных данных $\textbf{X}$ и выбирают релевантные признаки на основе (вневыборочной) производительности обучающегося алгоритма. Обучение линейной модели с различными комбинациями признаков и выбор признаков, дающих наилучшую вневыборочную производительность, было бы примером метода-обёртки.
	\item \textbf{Методы фильтрации} не используют обучающийся алгоритм на исходных данных $\textbf{X}$, а учитывают только статистические характеристики входных данных. Например, мы можем выбрать признаки, для которых корреляция между признаком и целевой переменной превышает пороговое значение корреляции.
	\item \textbf{Встроенные методы} — это собирательная группа техник, которые выполняют отбор признаков как часть процесса построения модели. LASSO является примером встроенного метода.
\end{enumerate}

Мы сосредоточимся на методах фильтрации, и в частности рассмотрим методы фильтрации, которые используют \textbf{взаимную информации}, для оценки того, какие признаки следует включить в уменьшенный набор данных $\textbf{X}_S$. Результирующий критерий приводит к NP-трудной задаче оптимизации, и мы обсудим несколько способов, которыми можно попытаться найти оптимальные решения этой задачи.

\subsubsection{Совместная взаимная информация как критерий для отбора признаков}

Предоположим следующее: в нашем распоряжении есть достаточно большой набор пар вида \enquote{целевая перменная --- вектор признаков}:
$$
	\mathcal{D} = \set{(y^{(i)}; \vx^{(i)})}_{i = 1}^{n},
$$
где $y^{(i)} \in \R, \vx^{(i)} \in \R^p$; Мы будем предполагать что пара $(y, \vx)$ это независимая реализация некоторой случайной величины $(\eta, \xi)$, $\eta$ --- со значениями в $\R$, $\xi$ со значениями в $\R^p$; Такая постановка позволяет говорить о взаимной информации $\I(\eta; \xi)$ между признаками $\xi$ и целевой переменной $\eta$;

Когда дело доходит до отбора признаков, мы хотим максимизировать взаимную информацию между подмножеством выбранных признаков $\xi_S$ и целевой переменной $\eta$:
$$
	\tilde{S} = \argmax_S \I(\xi_S; \eta), \quad \text{при условии } |S| = K, 
$$
где $K$ --- количество признаков, которые мы хотим выбрать. Эта величина называется \textbf{совместной взаимной информацией}, и её максимизация является NP-трудной задачей оптимизации, потому что множество возможных комбинаций признаков растёт экспоненциально.

С точки зрения данных, мы можем записать взаимную информацию следующим образом:
\[
	\mathcal{I} = \KL(p_{\xi_S, \eta}(x_S, y) | p_{\xi_S}(x_S) \cdot p_{\eta}(y)) = \E_{X_S \sim \xi_S, Y \sim \eta}[\log(\frac{p_{\xi_S, \eta}(X_S, Y)}{p_{\xi_S}(X_S) \cdot p_{\eta}(Y)})] \approx \frac{1}{n}\sum_{i = 1}^{n}\log(\frac{p_{\xi_S, \eta}(\vx_S^{(i)}, y^{(i)})}{p_{\xi_S}(\vx_S^{(i)}) \cdot p_{\eta}(y^{(i)})})
\]
В случае когда все данные дискретны, мы можем оценивать $p_{\xi_S, \eta}(X_S, Y)$  с помощью подсчета вхождений впар $(X_S, Y)$ в набор $\mathcal{D}$; В непрерывном случае все несколько усложняется, однако и там есть свои методы.

\subsubsection{Жадное решение задачи отбора признаков}

Самый простой подход к решению этой оптимизационной задачи --- жадный пошаговый алгоритм прямого отбора. Здесь признаки выбираются последовательно, по одному за раз. Пусть $S^{t - 1} = \{x_{f_1}, \ldots, x_{f_{t - 1}}\}$ --- множество выбранных признаков на шаге $t - 1$. Жадный метод выбирает следующий признак $f_t$ такой, что
$$
	f_t = \argmax_{\ r \notin S^{t - 1}} \I(\xi_{S^{t - 1} \cup r}; \eta). 
$$
Из правила цепочки мы можем записать:
$$
	\I(\xi_{S^{t-1} \cup r}; \eta) = \I(\xi_r; \eta | \xi_{S^{t - 1}}) + \underbrace{\sum_{k = 1}^{t-1} \I(\xi_{k}; \eta | \xi_{k - 1}, \ldots, \xi_1)}_{\const}
$$

Таким образом, жадный отбор признаков выбирает признаки, которые на каждом шаге дают наибольшее увеличение совместной взаимной информации. Однако даже такой подход, на каждом шаге требует огромное количество рассчетов. Теперь используем определение информации взаимодействия:
\[
	\I(\xi_r; \eta; \xi_{S^{t-1}}) = \I(\xi_r; \eta) - \I(\xi_r; \eta | \xi_{S^{t - 1}}) = \I(\xi_r; \xi_{S^{t - 1}}) - \I(\xi_r; \xi_{S^{t - 1}} | \eta)
\]
Тогда:
$$ 
f_t = \argmax_{\ r\notin S^{t - 1}} \underbrace{\I(\xi_r; \eta)}_{\text{релевантность}} - \underbrace{\bigl[\I(\xi_r; \xi_{S^{t - 1}}) - \I(\xi_r; \xi_{S^{t - 1}} | \eta) \bigr]}_{\text{избыточность}}.
$$

Это даёт нам понимание того, что по существу означает оптимизация данного критерия. Оптимизация этого критерия приводит к компромиссу между \textit{релевантностью} нового признака $x_i$ по отношению к цели $y$ и \textit{избыточностью} этой информации по сравнению с информацией, содержащейся в уже выбранных переменных $X_{S^{t - 1}}$.

\subsubsection{Аппроксимация меньшей размерности}

Даже с приведённым выше упрощением совместной взаимной информации, величины, включающие $S^{t - 1}$, всё ещё являются $(t - 1)$-мерными интегралами. Чтобы сделать задачу более разрешимой, многие подходы в литературе предлагают одну или обе из следующих аппроксимаций низкого порядка:

\begin{align}
I(x_i ; x_{S^{t - 1}}) &\approx \alpha \sum_{k = 1}^{t-1} I(x_{f_k}; x_i), \\
I(x_i ; x_{S^{t - 1}} | y) &\approx \beta \sum_{k = 1}^{t - 1} I(x_{f_k}; x_i | y).
\end{align}

Следовательно, задача оптимизации теперь упрощается до

$$
f_t = \arg\max_{i \notin S^{t - 1}} \underbrace{I(x_i; y)}_{\text{релевантность}} - \underbrace{\left[ \alpha \sum_{k = 1}^{t-1} I(x_{f_k}; x_i) - \beta \sum_{k = 1}^{t - 1} I(x_{f_k}; x_i | y) \right]}_{\text{избыточность}}
$$
где $\alpha$ и $\beta$ должны быть заданы.

\subsubsection{Семейство алгоритмов отбора признаков}

Эти параметры фактически определяют семейство критериев, основанных на взаимной информации, и мы можем получить некоторые известные примеры для конкретных значений $\alpha$ и $\beta$:

\begin{itemize}
	\item Совместная взаимная информация (JMI): $\alpha = \frac{1}{t - 1}$ и $\beta = \frac{1}{t - 1}$.
	\item Максимальная релевантность — минимальная избыточность (MRMR): $\alpha = \frac{1}{t - 1}$ и $\beta = 0$.
	\item Максимизация взаимной информации (MIM): $\alpha = 0$ и $\beta = 0$.
\end{itemize}

Различные значения $\alpha$ и $\beta$ имеют значение для того, что вы предполагаете о вероятностной структуре ваших данных. Они соответствуют одному (или обоим) из следующих предположений:

%\begin{quote}
	\textbf{Предположение 2}: Все признаки попарно условно независимы при условии класса, т.е.
	$$ p(x_i x_j | y) = p(x_i | y) p(x_j | y). $$
	Это означает, что $\sum I(X_j; X_k | y)$ будет равно нулю.
%\end{quote}

%\begin{quote}
	\textbf{Предположение 3}: Все признаки попарно независимы, т.е.
	$$ p(x_i x_j) = p(x_i) p(x_j). $$
	Это означает, что $\sum I(X_j; X_k)$ будет равно нулю.
%\end{quote}

Теперь мы можем видеть, что $\alpha$ и $\beta$ контролируют степень уверенности в одном из этих предположений:

\begin{itemize}
	\item Значение $\alpha$, близкое к нулю, указывает на более сильную уверенность в Предположении 3.
	\item Значение $\beta$, близкое к нулю, указывает на более сильную уверенность в Предположении 2.
	\item Все критерии отбора переменных принимают Предположение 1.
\end{itemize}


\subsection{Принцип максимума энтропии и логистическая регрессия}

Итак, допустим мы установили, что чем больше в данных информации о целевой переменной, тем лучше нижняя граница на вероятность ошибки. Неравенство Фано, однако, ничего не говорит о том, как же нам получить \enquote{хорошую} оценку, и вообще говоря не должно. Данные --- это одна история, предсказательная модель --- другая. Дизайн моделей, это отдельный большой и сложный раздел, который мы можем затронуть только поверхностно, рассмотрев один из принципов, которые можно использовать: принцип максимума энтропии.
	
\textbf{Принцип максимума энтропии.} Для описания набора данных $\mathcal{D}$, которые
\begin{enumerate}[compact]
	\item имеют вероятностную природу (т.е. $\mathcal{D}$ можно мыслить как реализацию некоторого случайного объекта, например, в самом простом случае $\mathcal{D} = \set{(y, \vx_i)}_{i = 1}^{n}, \vx_i \in \R^d$)
	\item распределение данных имеют некоторые ограничения (например, зафиксировано математическое ожидание $Y$;
\end{enumerate}





